\documentclass[11pt,oneside]{scrbook} % KOMA: pulito e flessibile

% ---------- LINGUA & CODIFICA ----------
\usepackage[italian]{babel}
\usepackage[T1]{fontenc}
\usepackage[utf8]{inputenc}

% ---------- TIPOGRAFIA CONSIGLIATA ----------
\usepackage{lmodern}     % font semplice e diffuso
\usepackage{microtype}   % micro-tipografia, migliore resa

% ---------- GEOMETRIA PAGINA ----------
% Scegli *una* delle due geometrie qui sotto.
% (1) Impostazioni consigliate per A4 (Italia) con ampio margine esterno:
\usepackage[
  a4paper,
  top=24mm,
  bottom=28mm,
  inner=24mm,   % margine interno (verso rilegatura)
  outer=30mm,   % margine esterno: ampio per le note a margine
  heightrounded
]{geometry}
\setlength{\headheight}{17pt} % evita sovrapposizioni con l'header
\setlength{\headsep}{6mm}     % distanza tra header e corpo

\usepackage{scrlayer-scrpage}
\clearpairofpagestyles
\ihead{\leftmark}
\ofoot*{\pagemark}
\automark{chapter}

% (2) Se vuoi replicare il formato del PDF caricato (US Letter 8.5x11"):
% \usepackage[
%   letterpaper,
%   top=25mm,
%   bottom=30mm,
%   inner=30mm,
%   outer=55mm,
%   heightrounded
% ]{geometry}

% ---------- STRUMENTI DI STRUTTURA ----------
\setcounter{secnumdepth}{3}
\setcounter{tocdepth}{3}

% ---------- STILE TITOLI IN STILE KOMA ----------
\setkomafont{chapter}{\normalfont\huge\bfseries}
\setkomafont{section}{\normalfont\Large\bfseries}
\setkomafont{subsection}{\normalfont\large\bfseries}
\setkomafont{subsubsection}{\normalfont\normalsize\bfseries}

% ---------- STILE FRONTESPIZIO ----------
\setkomafont{title}{\sffamily\bfseries\Huge}
\setkomafont{author}{\sffamily\large}
\setkomafont{publishers}{\sffamily\normalsize}
\setkomafont{date}{\large}

\RedeclareSectionCommand[beforeskip=1.5ex,afterskip=1.5ex]{chapter}
\RedeclareSectionCommand[beforeskip=1ex,afterskip=0.8ex]{section}
\RedeclareSectionCommand[beforeskip=0.8ex,afterskip=0.6ex]{subsection}
\RedeclareSectionCommand[beforeskip=0.6ex,afterskip=0.4ex]{subsubsection}

% ---------- AMBIENTI UTILI ----------
\usepackage{amsmath,amssymb,amsthm}
\usepackage{booktabs}
\usepackage{array}
\usepackage{graphicx}
\usepackage{hyperref}
\hypersetup{  % stile dei riferimenti
  colorlinks=true,
  linkcolor=black,
  citecolor=blue,
  urlcolor=blue
}
\usepackage{xcolor}
\usepackage{tcolorbox}
\usepackage{subcaption} 
\usepackage{upgreek} % \uptau

\usepackage{mathrsfs} % \mathscr{}

\usepackage{tikz}
\usetikzlibrary{intersections}
\usepackage{pgfplots} % grafici, creato in tikZ
\pgfplotsset{compat = 1.18}
\usetikzlibrary{
arrows.meta,
positioning,
calc,
decorations.markings,
angles,
quotes,
3d,
patterns
}

\usepackage{subcaption}

\usepackage[european]{circuitikz} % circuiti

\usepackage{comment} % \begin{comment}  \end{comment}

\usepackage{mathtools} % \coloneqq

\usepackage{siunitx}
\sisetup{separate-uncertainty=true}
%----------- TEOREMI E STRUTTURE VARIE (KOMA, AMSTHM)

% Stile "teorema" (corsivo)
\theoremstyle{plain}
\newtheorem{theorem}{Teorema}[chapter]
\newtheorem{proposition}[theorem]{Proposizione}
\newtheorem{lemma}[theorem]{Lemma}
\newtheorem{corollary}[theorem]{Corollario}

% Stile "definizione" (testo dritto)
\theoremstyle{definition}
\newtheorem{definition}[theorem]{Definizione}

% Stile "nota" (più discorsivo)
\theoremstyle{remark}
\newtheorem{remark}[theorem]{Osservazione}
\newtheorem{example}[theorem]{Esempio}


% ---------- NOTE A MARGINE (parole-chiave/sidenote) ----------
% \marginpar stampa nel margine esterno con oneside.
% Definiamo un comando comodo e graficamente coerente:
\newcommand{\mkeyword}[1]{%
  \marginpar{\raggedright\footnotesize\emph{#1}}%
}

% Se preferisci le note a sinistra, attiva la riga seguente:
% \reversemarginpar

% ---------- PARAGRAFI E SPAZI ----------
\setlength{\parindent}{0.8em}  % rientro tradizionale
\setlength{\parskip}{0.25em}   % piccolo spazio extra tra paragrafi

%----------- COMANDI AGGIUNTIVI -----------
% --- Comando per vettori colonna di lunghezza variabile ---

\newcommand{\colvec}[1]{%
  \begin{pmatrix}
    #1
  \end{pmatrix}%
}

% --- comando derivata --- parziale \der[ordine]{funzione}{variabile}
\newcommand{\der}[3][]{%
  \frac{\partial^{#1} #2}{\partial #3^{#1}}%
}

% --- comando versore --- 
\newcommand{\vers}[1]{
  \hat{\boldsymbol{#1}}
}

% --- epsilon ---
\newcommand{\eps}[0]{
  \varepsilon
}

% --- norma ---
\newcommand{\norm}[1]{
  \left\lVert #1 \right\rVert
}

% --- prodotto scalare ---
\newcommand{\sc}[1]{
  \langle #1 \rangle
}

% --- vettore ---
\renewcommand{\v}[1]{
  \mathbf{#1}
}

% ---------- BIBLIOGRAFIA ----------
\usepackage{csquotes} % \enquote{}
\usepackage[backend=biber,style=numeric,sorting=nyt]{biblatex}
\addbibresource{bibliografia.bib} % nome del file .bib

% ---------- FRONTESPIZIO ----------
\title{\textbf{\titolo}}
\author{Cosimo Pascarella}
\date{Anno Accademico 2025--2026}

\publishers{
  Email: \text{c.pascarella2@studenti.unipi.it}\\
  Versione: v0.1
}